% =====================================================================
%  Wspólna preambuła szablonów pism ZUS (świadczenia rodzicielskie 2026)
%  Silnik: XeLaTeX albo Tectonic (pdflatex NIE obsłuży polskich znaków w fontspec).
%  Kompilacja:  tectonic -X compile plik.tex
% =====================================================================
\documentclass[11pt,a4paper]{article}

\usepackage{fontspec}
\usepackage{polyglossia}
\setmainlanguage{polish}
\setotherlanguage{english}
% font domyślny (Latin Modern) — Tectonic go dostarcza

\usepackage[margin=2.5cm,a4paper]{geometry}
\usepackage{graphicx}
\usepackage{xcolor}
\usepackage{array}
\usepackage{enumitem}
\usepackage{fancyhdr}
\usepackage{hyperref}

\definecolor{zus}{HTML}{007A33}      % zieleń ZUS
\definecolor{szary}{HTML}{6B7280}
\definecolor{ramka}{HTML}{9CA3AF}

\hypersetup{colorlinks=true, linkcolor=zus, urlcolor=zus, pdftitle={Szablony pism ZUS}}

% --- pole do wypełnienia odręcznego -----------------------------------
% \Pole{nazwa} rysuje kropkowaną linię z podpisem pola pod spodem.
% Pole do wypełnienia: kropkowana linia, pod nią drobny podpis pola.
\newcommand{\Pole}[1]{%
  \textcolor{szary}{\dotfill}\,%
  {\scriptsize\textcolor{szary}{\textsc{#1}}}%
}

% --- blok warunkowy: treść pokazywana tylko, gdy pole prawdziwe --------
\newenvironment{warunek}{%
  \par\smallskip\noindent\begin{minipage}{\linewidth}%
  \color{black}%
}{%
  \end{minipage}\par\smallskip%
}

% --- nagłówek pisma ---------------------------------------------------
\newcommand{\NaglowekPisma}[1]{%
  \par\vspace{0.8em}%
  \begin{center}%
    {\large\bfseries\MakeUppercase{#1}}%
  \end{center}%
  \vspace{0.4em}%
}

% --- blok nadawcy (pierwsze linie pisma) -------------------------------
\newenvironment{nadawca}{%
  \begin{flushleft}\small
}{%
  \end{flushleft}\normalsize
}

% --- stopka z numerem strony -----------------------------------------
\pagestyle{fancy}
\fancyhf{}
\renewcommand{\headrulewidth}{0pt}
\fancyfoot[C]{\textcolor{szary}{\small\thepage}}

% --- metadane szablonu na początku pisma ------------------------------
\newcommand{\Metadane}[3]{%
  \begin{flushright}\footnotesize\textcolor{szary}%
    \textbf{#1}\\[2pt] {\itshape #2}\\[2pt] #3%
  \end{flushright}\vspace{0.5em}%
}

\setlength{\parindent}{0pt}
\setlength{\parskip}{0.6em}

\begin{document}
\section*{Decyzje ZUS --- struktura pisma}

{\small\color{szary}Decyzji ZUS nie generujemy. Poniższa struktura służy zrozumieniu pisma i przygotowaniu odwołania.}\par\medskip

\subsection*{Decyzja przyznająca}

\begin{itemize}\item Nagłówek: Zakład Ubezpieczeń Społecznych, Oddział w \{\{oddzial\}\}\end{itemize}

\begin{itemize}\item Znak sprawy / sygnatura: \{\{znak\_sprawy\}\}\end{itemize}

\begin{itemize}\item Nazwa i adres ubezpieczonego\end{itemize}

\begin{itemize}\item Rozstrzygnięcie: przysługuje / przyznaję świadczenie \{\{benefit\}\} za okres od \{\{od\}\} do \{\{do\}\}\end{itemize}

\begin{itemize}\item Podstawa prawna rozstrzygnięcia: \{\{podstawa\}\}\end{itemize}

\begin{itemize}\item Uzasadnienie faktyczne: stan faktyczny sprawy\end{itemize}

\begin{itemize}\item Uzasadnienie prawne: art. \{\{art\}\} ustawy zasiłkowej\end{itemize}

\begin{itemize}\item Kwota zasiłku: \{\{kwota\}\} zł\end{itemize}

\begin{itemize}\item Pouczenie o terminie odwołania (14 dni, art. 63 ustawy zasiłkowej do ZUS, potem do sądu pracy)\end{itemize}

\begin{itemize}\item Pieczęć i podpis \{\{dystrybutor\}\}\end{itemize}

{\footnotesize\color{szary}\itshape DECYZJA POZOSTAJE NARZĘDZIEM ZUS --- nie generujemy jej za użytkownika.
Struktura służy do zrozumienia pisma i przygotowania odwołania.
}

\subsection*{Odmowa / decyzja odmowna}

\begin{itemize}\item Rozstrzygnięcie: odmawiam przyznania \{\{benefit\}\}\end{itemize}

\begin{itemize}\item Uzasadnienie faktyczne: \{\{stan\_faktyczny\}\}\end{itemize}

\begin{itemize}\item Uzasadnienie prawne: art. \{\{art\}\} ustawy zasiłkowej\end{itemize}

\begin{itemize}\item Pouczenie o odwołaniu (14 dni do ZUS)\end{itemize}

\textbf{Wzór odwołania}\par

\NaglowekPisma{ODWOŁANIE OD DECYZJI ZAKŁADU UBEZPIECZEŃ SPOŁECZNYCH}

Na podstawie art. 63 ustawy zasiłkowej odwołuję się od decyzji ZUS
z dnia \Pole{data decyzji}, znak \Pole{znak sprawy}, odmawiającej mi prawa do \Pole{benefit}.

Uzasadnienie:
\Pole{uzasadnienie odwolania}

\Pole{data pisma}
.....................................
(podpis)
\end{document}
